\documentclass[11pt,a4paper]{article}

% ---------------------------------------------------------
% PREAMBLE (stable, warning-free, Overleaf-safe)
% ---------------------------------------------------------
\usepackage[utf8]{inputenc}
\usepackage[T1]{fontenc}
\usepackage{lmodern}
\usepackage{amsmath, amssymb, amsfonts}
\usepackage{bm}
\usepackage{geometry}
\usepackage{graphicx}
\usepackage{hyperref}
\usepackage{cite}
\usepackage{authblk}
\usepackage{physics}
\usepackage{mathtools}

\geometry{margin=2.4cm}

\title{\textbf{Companion XLII: The Kinematic Sunyaev--Zel'dovich Effect in the Relaxation-Driven Cyclic Cosmology}}
\author{Michael Lehmann \\ \texttt{mi.lehmann@gmx.de}}
\date{April 2026}

\begin{document}
\maketitle

\begin{abstract}
The kinematic Sunyaev--Zel'dovich (kSZ) effect provides a direct probe of cosmic 
velocity fields and the dynamical evolution of large-scale structure. In the 
standard cosmological model, the kSZ signal arises from the Doppler shift of CMB 
photons scattered by free electrons moving with bulk velocities. In the 
Relaxation-Driven Cyclic Cosmology (RDCC), the scale-dependent evolution of the 
gravitational potential modifies the growth of velocity fields, the coherence of 
bulk flows, and the coupling between baryons and dark matter. This Companion 
develops a continuous description of the kSZ effect in RDCC, deriving how the 
relaxation scale influences the velocity divergence, the pairwise momentum field, 
and the angular power spectrum of the kSZ signal. The resulting framework provides 
a new observational window into the dynamical structure of the RDCC Universe.
\end{abstract}
% ---------------------------------------------------------
\section{Introduction}

The kinematic Sunyaev--Zel'dovich effect is a Doppler shift imprinted on the cosmic 
microwave background by free electrons moving with bulk velocities. Unlike the 
thermal SZ effect, which probes the thermal energy of the intracluster medium, the 
kSZ effect is sensitive to the momentum field of baryons and thus to the dynamical 
evolution of large-scale structure.

In the standard cosmological model, the velocity field is governed by the linear 
growth rate and the continuity equation. In RDCC, the gravitational potential 
evolves differently. The relaxation mechanism suppresses the decay of small-scale 
modes and enhances the coherence of large-scale modes. This modifies the velocity 
divergence, the amplitude of bulk flows, and the coupling between baryons and dark 
matter. The kSZ effect becomes a direct tracer of the relaxation scale 
$k_{\mathrm{rel}}$ and the temporal structure of the gravitational field.
% ---------------------------------------------------------
\section{Velocity Fields in RDCC}

The cosmic velocity field is governed by the Euler equation,
\begin{equation}
    \frac{\partial \mathbf{v}}{\partial a}
    + \frac{1}{a} (\mathbf{v} \cdot \nabla)\mathbf{v}
    = -\frac{1}{aH} \nabla \Phi_{\mathrm{RDCC}}.
\end{equation}

In RDCC, the gravitational potential takes the form
\begin{equation}
    \Phi_{\mathrm{RDCC}}(k,a)
    = \Phi(k,a)
      \left[
        1 - \chi_{\Phi}
        \left(\frac{k}{k_{\mathrm{rel}}}\right)^{\alpha}
      \right].
\end{equation}

The velocity divergence,
\begin{equation}
    \theta = \nabla \cdot \mathbf{v},
\end{equation}
is therefore modified by the relaxation mechanism. Small-scale modes contribute 
less to the divergence, while large-scale modes retain their coherence. The 
velocity field becomes smoother and more coherent on scales near $k_{\mathrm{rel}}$.
% ---------------------------------------------------------
\section{The RDCC kSZ Signal}

The kSZ temperature shift is given by
\begin{equation}
    \frac{\Delta T_{\mathrm{kSZ}}}{T_{\gamma}}
    = -\sigma_{T}
      \int n_{e}(\mathbf{x})\,
      \frac{\mathbf{v} \cdot \hat{\mathbf{n}}}{c}\,
      dl,
\end{equation}
where $\hat{\mathbf{n}}$ is the line-of-sight direction.

In RDCC, both the electron density and the velocity field carry the imprint of the 
relaxation scale. The kSZ signal becomes
\begin{equation}
    \Delta T_{\mathrm{kSZ,RDCC}}
    = \Delta T_{\mathrm{kSZ}}
      \left[
        1 - \chi_{\mathrm{kSZ}}
        \left(\frac{k}{k_{\mathrm{rel}}}\right)^{\alpha}
      \right].
\end{equation}

This modification suppresses small-scale kSZ fluctuations and enhances large-scale 
coherence.
% ---------------------------------------------------------
\section{Pairwise Momentum and Bulk Flows}

The pairwise momentum field,
\begin{equation}
    p_{12}(r)
    = \left\langle
        \left[
            1 + \delta_{e}(\mathbf{x})
        \right]
        \mathbf{v}(\mathbf{x})
        \cdot \hat{\mathbf{r}}
    \right\rangle,
\end{equation}
is a key observable for the kSZ effect.

In RDCC, the suppression of small-scale potential evolution reduces the amplitude 
of small-scale pairwise velocities, while the enhanced coherence of large-scale 
modes increases the amplitude of bulk flows. The pairwise momentum becomes
\begin{equation}
    p_{12,\mathrm{RDCC}}(r)
    = p_{12}(r)
      \left[
        1 - \chi_{p}
        \left(\frac{1}{k_{\mathrm{rel}} r}\right)^{\alpha}
      \right].
\end{equation}

This produces a characteristic scale-dependent modulation that reflects the 
relaxation scale.
% ---------------------------------------------------------
\section{The kSZ Power Spectrum}

The angular power spectrum of the kSZ effect is given by
\begin{equation}
    C_{\ell}^{\mathrm{kSZ}}
    = \int dk\,
      P_{q}(k)\,
      j_{\ell}^{2}(kr),
\end{equation}
where $P_{q}(k)$ is the power spectrum of the momentum field.

In RDCC, the momentum power spectrum becomes
\begin{equation}
    P_{q,\mathrm{RDCC}}(k)
    = P_{q}(k)
      \left[
        1 - \chi_{q}
        \left(\frac{k}{k_{\mathrm{rel}}}\right)^{\alpha}
      \right].
\end{equation}

This modification suppresses small-scale kSZ power and enhances large-scale 
coherence, producing a distinctive RDCC signature in the kSZ power spectrum.
% ---------------------------------------------------------
\section{Conclusion}

The kinematic Sunyaev--Zel'dovich effect in the Relaxation-Driven Cyclic Cosmology 
reflects the scale-dependent evolution of the gravitational potential. The 
relaxation mechanism modifies the velocity divergence, the amplitude of bulk flows, 
and the coupling between baryons and dark matter. These effects produce distinctive 
signatures in the kSZ temperature shift, the pairwise momentum field, and the kSZ 
power spectrum. The present Companion establishes the theoretical foundation for 
future studies of cosmic velocity fields, baryon dynamics, and the interplay 
between matter and radiation in RDCC.

% ========================
% References
% ========================

\begin{thebibliography}{10}

\bibitem{Flagship} Lehmann, M. (2026). Relaxation-Driven Cyclic Cosmology (RDCC): A Minimal CPT-Symmetric Framework with a Single Parameter. Flagship Paper. Zenodo: \url{https://zenodo.org/records/19744958} (v43.0 or later).

\bibitem{Zenodo} The complete RDCC companion series (I--LVII) is available at the same Zenodo record above.

% Thematisch relevante Companions
\bibitem{CompXVI} Companion XVI: Boltzmann Code and Modified Hierarchy.
\bibitem{CompXXXI} Companion XXXI: RDCC Halo Model and Non-Linear Structure.
\bibitem{CompXXXII} Companion XXXII--XXXIX: Galaxy--Halo Connection, Cosmic Web, Lensing, RSD, ISW, tSZ, Rotation Curves, CGM.

% Wichtige externe Referenzen
\bibitem{Planck2020} Planck Collaboration (2020), Astron. Astrophys. 641, A6.
\bibitem{DESI2024} DESI Collaboration (2024/2025), arXiv: [aktuelle $f\sigma_8$/BAO-Paper].
\bibitem{JWST2024} Maiolino et al. (2024), Nature (JWST high-$z$ galaxies/quasars).
\bibitem{Kaiser1987} Kaiser, N. (1987), Mon. Not. R. Astron. Soc. 227, 1 (RSD).
\bibitem{Bartelmann2001} Bartelmann, M. \& Schneider, P. (2001), Phys. Rep. 340, 291 (Weak Lensing Review).
\bibitem{HuOkamoto2002} Hu, W. \& Okamoto, T. (2002), Astrophys. J. 574, 566 (CMB Lensing).
\bibitem{LewisChallinor2006} Lewis, A. \& Challinor, A. (2006), Phys. Rep. 429, 1 (CMB Lensing Review).
\bibitem{NFW1997} Navarro, J. F., Frenk, C. S. \& White, S. D. M. (1997), Astrophys. J. 490, 493 (Halo Profiles).

\end{thebibliography}

\end{document}
